% Generated draft: requires mathematical and source review.
\documentclass[11pt]{article}
\usepackage[margin=25mm]{geometry}
\usepackage{fontspec}
\setmainfont{Latin Modern Roman}
\usepackage{amsmath,amssymb,mathtools,unicode-math}
\setmathfont{Latin Modern Math}
\usepackage{enumitem,needspace,xurl,hyperref}
\setlength{\parindent}{0pt}
\setlength{\parskip}{5pt}
\setlength{\emergencystretch}{3em}
\newcommand{\N}{\mathbb{N}}
\newcommand{\Z}{\mathbb{Z}}
\newcommand{\Q}{\mathbb{Q}}
\newcommand{\R}{\mathbb{R}}
\newcommand{\C}{\mathbb{C}}
\newcommand{\F}{\mathbb{F}}
\newcommand{\A}{\mathbb{A}}
\newcommand{\PP}{\mathbb{P}}
\newcommand{\E}{\mathbb{E}}
\newcommand{\eps}{\varepsilon}
\newcommand{\dd}{\,\mathrm{d}}
\begin{document}
\section*{UnsolvedMath ID: 800872}
\textbf{Source problem code:} catalogue entry
\subsection{Definitions and mathematical statement}
The source leaves the ambient domain and conventions implicit. The following is the usual domain formulation: let \(n\geq 1\) be an integer and let \(\Omega\subset\C^n\) be a nonempty domain. Write \(\operatorname{PSH}(\Omega)\) for the upper semicontinuous functions \(u:\Omega\to[-\infty,\infty)\) whose restrictions to the connected components of intersections with complex lines are subharmonic or identically \(-\infty\), excluding the function identically \(-\infty\) on \(\Omega\). Values \(-\infty\) at individual points are allowed; no local boundedness or smoothness hypothesis is imposed.

A function \(u\in\operatorname{PSH}(\Omega)\) is called \emph{maximal on \(\Omega\)} if, for every open set \(G\Subset\Omega\) and every upper semicontinuous function \(v:\overline G\to[-\infty,\infty)\) that is plurisubharmonic on \(G\),
\[
 v\leq u\ \text{on }\partial G
 \quad\Longrightarrow\quad
 v\leq u\ \text{on }G.
\]
Here \(G\Subset\Omega\) means that \(\overline G\) is compact and contained in \(\Omega\). Call \(u\) \emph{locally maximal} if every \(x\in\Omega\) has an open neighborhood \(U_x\subset\Omega\) on which \(u|_{U_x}\) is maximal in this comparison sense.

\textbf{Question 3 (attributed in the source to Blocki).} For every such \(n,\Omega\), and \(u\in\operatorname{PSH}(\Omega)\), is
\[
 u\text{ maximal on }\Omega
 \quad\Longleftrightarrow\quad
 \forall x\in\Omega\ \exists\text{ an open neighborhood }U_x\subset\Omega
 \text{ such that }u|_{U_x}\text{ is maximal on }U_x?
\]

The supplied background states that in complex dimension one maximal plurisubharmonic functions are harmonic and hence smooth, whereas in higher dimensions maximal functions can be irregular; it gives functions depending on only \(n-1\) complex variables as examples. It also reports, with reference label [Blo04], that
\[
 \varphi\in DMA(\Omega)
 \quad\Longrightarrow\quad
 \bigl(\varphi\text{ maximal}\ \Longleftrightarrow\ MA(\varphi)=0\bigr),
\]
and consequently that maximality is local within \(DMA(\Omega)\). Here \(MA\) denotes the complex Monge--Amp\`ere operator, and \(DMA(\Omega)\) denotes its domain of definition as used by the source. The supplied record does not specify the precise construction of this domain. These are supplied contextual assertions, not additional hypotheses on the question. The question concerns arbitrary plurisubharmonic functions, including those outside that domain. The domain convention and the exclusion of the identically \(-\infty\) function above require confirmation against the original source.
\subsection{Short English statement}
Must a plurisubharmonic function that satisfies the maximality comparison principle near every point satisfy it on the entire domain? The supplied source attributes this question to Blocki and reports locality within the domain of definition of the complex Monge--Amp\`ere operator. Its registry status is \emph{open}; this draft does not independently verify that status.
\subsection{Sources}
\begin{itemize}
\item Local UnsolvedMath registry, permanent ID \(800872\), source problem code \texttt{catalogue entry}, title \emph{Open Problems in Pluripotential Theory --- Question 3}. The supplied record lists Dinew, Guedj, and Zeriahi as proposers, with proposed year 2015, and labels the problem open.
\item Dinew, Guedj, and Zeriahi, \emph{Open Problems in Pluripotential Theory} (2015), Question 3: \url{https://arxiv.org/abs/1511.00705}. The supplied provenance records access on 2026-10-16 and extraction by \texttt{arxiv\_latex}; these are supplied provenance details, not a claim of independent access or verification.
\item The question is individually attributed to Blocki in the supplied excerpt. The background cites Blocki with reference label [Blo04] for the characterization of maximality by vanishing Monge--Amp\`ere measure in \(DMA(\Omega)\). No full bibliographic entry for [Blo04] is supplied.
\end{itemize}
\end{document}
